\documentclass[CJK,13pt]{beamer}
\usepackage{CJKutf8}
\usepackage{beamerthemesplit}
\usetheme{Warsaw}
\useoutertheme[footline=authortitle]{miniframes}
\usepackage{amsmath}
\usepackage{amssymb}
\usepackage{graphicx}
\usepackage{eufrak}
\usepackage{color}
\usepackage{slashed}
\usepackage{simplewick}
\usepackage{tikz}
\usepackage{tcolorbox}
\usepackage{ulem}
%%figures
\def\lfig#1#2{\includegraphics[width=#1 in]{#2}}
\def\tfig#1#2{\includegraphics[height=#1 in]{#2}}
\def\addfig#1#2{\begin{center}\includegraphics[width=#1 in]{#2}\end{center}}
\def\question{\includegraphics[width=0.3in]{why.jpg}\,}
\def\answer{\includegraphics[width=0.3in]{baozou_haha.png}\,}
\def\wulian{\includegraphics[width=0.18in]{emoji_wulian.jpg}}
\def\bigwulian{\includegraphics[width=0.35in]{emoji_wulian.jpg}}
\def\bye{\includegraphics[width=0.18in]{emoji_bye.jpg}}
\def\bigbye{\includegraphics[width=0.35in]{emoji_bye.jpg}}
\def\huaixiao{\includegraphics[width=0.18in]{emoji_huaixiao.jpg}}
\def\bighuaixiao{\includegraphics[width=0.35in]{emoji_huaixiao.jpg}}
\def\jianxiao{\includegraphics[width=0.18in]{emoji_jianxiao.jpg}}
\def\bigjianxiao{\includegraphics[width=0.35in]{emoji_jianxiao.jpg}}
\def\haoqi{\includegraphics[width=0.18in]{emoji_haoqi.jpg}}
%% colors
\def\blacktext#1{{\color{black}#1}}
\def\bluetext#1{{\color{blue}#1}}
\def\redtext#1{{\color{red}#1}}
\def\darkbluetext#1{{\color[rgb]{0,0.2,0.6}#1}}
\def\skybluetext#1{{\color[rgb]{0.2,0.7,1.}#1}}
\def\cyantext#1{{\color[rgb]{0.,0.5,0.5}#1}}
\def\greentext#1{{\color[rgb]{0,0.7,0.1}#1}}
\def\darkgray{\color[rgb]{0.2,0.2,0.2}}
\def\lightgray{\color[rgb]{0.6,0.6,0.6}}
\def\gray{\color[rgb]{0.4,0.4,0.4}}
\def\blue{\color{blue}}
\def\red{\color{red}}
\def\orange{\color[rgb]{1.,0.8,0.}}
\def\green{\color{green}}
\def\darkgreen{\color[rgb]{0,0.4,0.1}}
\def\darkblue{\color[rgb]{0,0.2,0.6}}
\def\skyblue{\color[rgb]{0.2,0.7,1.}}
%%control
\def\ulim#1{\underset{#1}{\lim}}
\def\heaviside{\,\mathrm{h}\,}
\def\bral{\left(\begin{array}{l}}
\def\brar{\end{array}\right)}
\def\brall{\left(\begin{array}{ll}}
\def\brarr{\end{array}\right)}
\def\bralll{\left(\begin{array}{lll}}
\def\brarrr{\end{array}\right)}
\def\branchl{\left\{\begin{array}{l}}
\def\branchr{\end{array}\right.}
\def\branchll{\left\{\begin{array}{ll}}
\def\branchrr{\end{array}\right.}
\def\branchlll{\left\{\begin{array}{lll}}
\def\branchrrr{\end{array}\right.}
\def\sfgamma#1{\,\Gamma\left( #1 \right)\,}
\def\be{\begin{equation}}
\def\ee{\nonumber\end{equation}}
\def\bea{\begin{eqnarray}}
\def\eea{\nonumber\end{eqnarray}}
\def\bch{\begin{CJK}{UTF8}{gbsn}}
\def\ech{\end{CJK}}
\def\bitem{\begin{itemize}}
\def\eitem{\end{itemize}}
\def\bcenter{\begin{center}}
\def\ecenter{\end{center}}
\def\bex{\begin{minipage}{0.2\textwidth}\includegraphics[width=0.6in]{jugelizi.png}\end{minipage}\begin{minipage}{0.76\textwidth}}
\def\eex{\end{minipage}}
\def\chtitle#1{\frametitle{\bch#1\ech}}
\def\bmat#1{\left(\begin{array}{#1}}
\def\emat{\end{array}\right)}
\def\bcase#1{\left\{\begin{array}{#1}}
\def\ecase{\end{array}\right.}
\def\bmini#1{\begin{minipage}{#1\textwidth}}
\def\emini{\end{minipage}}
\def\tbox#1{\begin{tcolorbox}#1\end{tcolorbox}}
\def\pfrac#1#2#3{\left(\frac{\partial #1}{\partial #2}\right)_{#3}}
\def\res#1#2{\,\mathrm{res}\,#1\left(#2\right)\,}
\def\newt#1#2{\left(\begin{array}{c}#1\\ #2\end{array}\right)}
\def\reof#1{\,\mathrm{Re}{\left(#1\right)\,}}
\def\imof#1{\,\mathrm{Im}{\left(#1\right)\,}}
\def\Arg#1{\,\mathrm{Arg}\,#1\,}
%%symbols
\def\intfull{\int_{-\infty}^\infty}
\def\inthalf{\int_0^\infty}
\def\bropt{\,(\ \ \ )}
\def\sech{\mathrm{sech}\,}
\def\csch{\mathrm{csch}\,}
\def\sone{$\star$}
\def\stwo{$\star\star$}
\def\sthree{$\star\star\star$}
\def\sfour{$\star\star\star\star$}
\def\sfive{$\star\star\star\star\star$}
\def\rint{{\int_\leftrightarrow}}
\def\roint{{\oint_\leftrightarrow}}
\def\stdHf{{\textit{\r H}_f}}
\def\deltaH{{\Delta \textit{\r H}}}
\def\ii{{\dot{\imath}}}
\def\skipline{{\vskip0.1in}}
\def\skiplines{{\vskip0.2in}}
\def\lagr{{\mathcal{L}}}
\def\hamil{{\mathcal{H}}}
\def\vecv{{\mathbf{v}}}
\def\vecf{{\mathbf{f}}}
\def\vecg{{\mathbf{g}}}
\def\vecupsilon{{\mathbf{\upsilon}}}
\def\vecu{{\mathbf{u}}}
\def\vecj{{\mathbf{j}}}
\def\vecx{{\mathbf{x}}}
\def\vecy{{\mathbf{y}}}
\def\vecz{{\mathbf{z}}}
\def\veck{{\mathbf{k}}}
\def\vecp{{\mathbf{p}}}
\def\vecn{{\mathbf{n}}}
\def\vecA{{\mathbf{A}}}
\def\vecP{{\mathbf{P}}}
\def\vecsigma{{\mathbf{\sigma}}}
\def\hatJn{{\hat{J_\vecn}}}
\def\hatJx{{\hat{J_x}}}
\def\hatJy{{\hat{J_y}}}
\def\hatJz{{\hat{J_z}}}
\def\hatj#1{\hat{J_{#1}}}
\def\hatphi{{\hat{\phi}}}
\def\hatq{{\hat{q}}}
\def\hatpi{{\hat{\pi}}}
\def\vel{\upsilon}
\def\Dint{{\mathcal{D}}}
\def\adag{{\hat{a}^\dagger}}
\def\bdag{{\hat{b}^\dagger}}
\def\cdag{{\hat{c}^\dagger}}
\def\ddag{{\hat{d}^\dagger}}
\def\hata{{\hat{a}}}
\def\hatb{{\hat{b}}}
\def\hatc{{\hat{c}}}
\def\hatd{{\hat{d}}}
\def\hatD{{\,\hat{D}}}
\def\hatN{{\hat{N}}}
\def\hatH{{\hat{H}}}
\def\hatp{{\hat{p}}}
\def\Fup{{F^{\mu\nu}}}
\def\Fdown{{F_{\mu\nu}}}
\def\newl{\nonumber \\}
\def\vece{\mathrm{e}}
\def\calM{{\mathcal{M}}}
\def\calT{{\mathcal{T}}}
\def\calR{{\mathcal{R}}}
\def\barpsi{\bar{\psi}}
\def\baru{\bar{u}}
\def\barv{\bar{\upsilon}}
\def\qeq{\stackrel{?}{=}}
\def\ftf{\stackrel{\mathcal{FT}}{\Longrightarrow}}
\def\ftb{\stackrel{\mathcal{FT}}{\Longleftarrow}}
\def\ftfb{\stackrel{\mathcal{FT}}{\Longleftrightarrow}}
\def\ltf{\stackrel{\mathcal{LT}}{\Longrightarrow}}
\def\ltb{\stackrel{\mathcal{LT}}{\Longleftarrow}}
\def\ltfb{\stackrel{\mathcal{LT}}{\Longleftrightarrow}}
\def\torder#1{\mathcal{T}\left(#1\right)}
\def\rorder#1{\mathcal{R}\left(#1\right)}
\def\contr#1#2{\contraction{}{#1}{}{#2}#1#2}
\def\trof#1{\mathrm{Tr}\left(#1\right)}
\def\trace{\mathrm{Tr}}
\def\comm#1{\ \ \ \left(\mathrm{used}\ #1\right)}
\def\tcomm#1{\ \ \ (\text{#1})}
\def\slp{\slashed{p}}
\def\slk{\slashed{k}}
\def\calp{{\mathfrak{p}}}
\def\veccalp{\mathbf{\mathfrak{p}}}
\def\Tthree{T_{\tiny \textcircled{3}}}
\def\pthree{p_{\tiny \textcircled{3}}}
\def\dbar{{\,\mathchar'26\mkern-12mu d}}
\def\erf{\mathrm{erf}}
\def\const{\mathrm{const.}}
\def\pheat{\pfrac p{\ln T}V}
\def\vheat{\pfrac V{\ln T}p}
%%units
\def\fdeg{{^\circ \mathrm{F}}}
\def\cdeg{^\circ \mathrm{C}}
\def\atm{\,\mathrm{atm}}
\def\angstrom{\,\text{\AA}}
\def\SIL{\,\mathrm{L}}
\def\SIkm{\,\mathrm{km}}
\def\SIyr{\,\mathrm{yr}}
\def\SIGyr{\,\mathrm{Gyr}}
\def\SIV{\,\mathrm{V}}
\def\SImV{\,\mathrm{mV}}
\def\SIeV{\,\mathrm{eV}}
\def\SIkeV{\,\mathrm{keV}}
\def\SIMeV{\,\mathrm{MeV}}
\def\SIGeV{\,\mathrm{GeV}}
\def\SIcal{\,\mathrm{cal}}
\def\SIkcal{\,\mathrm{kcal}}
\def\SImol{\,\mathrm{mol}}
\def\SIN{\,\mathrm{N}}
\def\SIHz{\,\mathrm{Hz}}
\def\SIm{\,\mathrm{m}}
\def\SIcm{\,\mathrm{cm}}
\def\SIfm{\,\mathrm{fm}}
\def\SImm{\,\mathrm{mm}}
\def\SInm{\,\mathrm{nm}}
\def\SImum{\,\mathrm{\mu m}}
\def\SIJ{\,\mathrm{J}}
\def\SIW{\,\mathrm{W}}
\def\SIkJ{\,\mathrm{kJ}}
\def\SIs{\,\mathrm{s}}
\def\SIkg{\,\mathrm{kg}}
\def\SIg{\,\mathrm{g}}
\def\SIK{\,\mathrm{K}}
\def\SImmHg{\,\mathrm{mmHg}}
\def\SIPa{\,\mathrm{Pa}}
%page
\def\secpage#1#2{\begin{frame}\bcenter{\bf \Huge #1} \ecenter \skipline \tbox{\bcenter #2 \ecenter}\end{frame}}
\def\append#1#2{\secpage{附录#1}{#2}}
\def\thinka#1{\begin{frame}\frametitle{思考题}\bcenter\lfig{0.4}{think.jpg}  \skipline #1 \ecenter\end{frame}}
\def\thinkb#1{\begin{frame}\frametitle{思考题}\bcenter\lfig{0.5}{think1.jpg}  \skipline #1 \ecenter\end{frame}}
\def\thinkc#1{\begin{frame}\frametitle{思考题}\bcenter\lfig{0.5}{think2.jpg}  \skipline #1 \ecenter\end{frame}}
\def\thinkd#1{\begin{frame}\frametitle{思考题}\bcenter\lfig{0.5}{think3.jpg}  \skipline #1 \ecenter\end{frame}}
\def\thinke#1{\begin{frame}\frametitle{思考题}\bcenter\lfig{0.5}{think4.jpg}  \skipline #1 \ecenter\end{frame}}

\def\courseurl{http://zhiqihuang.top/cosm}

\def\tpage#1#2{
\title{COSM \S{#1}  #2}
  \author{Zhiqi Huang}

\begin{frame}
\begin{center}
{\bf \Large Cosmology}

{\vskip 0.1in}

{\Large \S #1 #2}

{\vskip 0.2in}

{Lecturer: 黄志琦}

\vskip 0.2in

\courseurl

\end{center}
\end{frame}
}


  \date{}
  \begin{document}
  \bch
  \tpage{6}{Extended Cosmological Models}

  \section{wCDM}
  \secpage{$w$CDM模型}{$$\rho_{\rm DE} \propto a^{-3(1+w)}$$}

  \begin{frame}
    \frametitle{宇宙学常数面临的疑难}
    宇宙学常数(真空能) $\Lambda$ 作为暗能量的标准解释，看起来既高大上又对得上绝大多数观测数据，美滋滋。但其实，它至少面临着两个理论难题：
    \bitem
  \item{{\blue fine tuning problem:} 按照量子场论预言的真空能至少也得是 $\mathrm{TeV}^4$ 的量级，和实际观测到的 $\mathrm{meV}^4$ 差了那么亿点点。}
  \item{{\blue coincidence problem:} 按照 $\Lambda$CDM 模型，在宇宙演化历史中 $\rho_m$ 至少变化了几十个数量级，而 $\rho_\Lambda$ 不变。那为什么目前 $\rho_m$ 和 $\rho_\Lambda$ 是同一数量级的呢？}
    \eitem
  \end{frame}

  \begin{frame}
    \frametitle{$w$CDM}
    既然 $\Lambda$CDM 未必稳，那么就搞亿点点事情？

    \addfig{1.7}{gaoshiqing.jpg}

    最不废脑子的办法就是假设暗能量是状态方程（压强和密度之比）为常数 $w$ 的，和其他成分仅有引力相互作用的未知流体。这个模型称为{\blue $w$CDM模型}。
  \end{frame}

  \begin{frame}
    \frametitle{$w$CDM 的基本计算}
    请先用能量守恒证明 $w$CDM 模型里的暗能量密度：
    $$\rho_{\rm DE} \propto a^{-3(1+w)}$$
    然后证明
   {\small $$H(z)=H_0\sqrt{\Omega_r(1+z)^4+\Omega_m(1+z)^3 + \Omega_k(1+z)^2+\Omega_{\rm DE}(1+z)^{3(1+w)}}$$}

    \skiplines
    
    （各种距离，共动体积等的表达式，只要替换掉 $H(z)$，其余全都不用变。）
  \end{frame}

  \section{CPL model}
  \secpage{$w_0$-$w_a$ 模型}{$$\rho_{\rm DE} \propto e^{-3\left[(1+w_0+w_a)\ln a + w_a(1-a)\right]}$$}
  
  \begin{frame}
    \frametitle{$w_0$-$w_a$ 模型}
    \lfig{0.2}{why.jpg}等等，刚才你为什么假设暗能量的状态方程是一个常数？
    
    \lfig{0.2}{baozou_haha.png} 为了省脑子……

    \skiplines
    
    除了常数之外，最不废脑子（这是重点）的就是线性函数。在 $w_0$-$w_a$模型（有时也按三位提出者Chevallier-Polarski-Linder命名为CPL模型) 里，假设暗能量的状态方程为
    $$w_0+w_a(1-a)$$
      这里的 $w_0, w_a$ 为常数， $a$ 是尺度因子。

      \skipline
      
      （显然，$w$CDM 模型也可以看成 CPL模型里 $w_a=0$ 的特殊情形。）
  \end{frame}

  \begin{frame}
    \frametitle{$w_0$-$w_a$ 模型的基本计算}
    请先用能量守恒证明 $w_0$-$w_a$ 模型里的暗能量密度：
    $$\rho_{\rm DE} \propto e^{-3\left[(1+w_0+w_a)\ln a + w_a(1-a)\right]}$$
    然后证明
   {\scriptsize $$H(z)=H_0\sqrt{\Omega_r(1+z)^4+\Omega_m(1+z)^3 + \Omega_k(1+z)^2+\Omega_{\rm DE}e^{3\left[(1+w_0+w_a)\ln (1+z) - w_a\frac{z}{1+z}\right]}}$$}

    \skiplines
    
    （各种距离，共动体积等的表达式，只要替换掉 $H(z)$，其余全都不用变。）
  \end{frame}

  \section{Cosmic Neutrino Background}

  \secpage{Cosmic Neutrino Background}{刚才省下的亿点点脑细胞都得废在这里了！}

  \begin{frame}
    \frametitle{中微子、正负电子和光子的一段往事}
    大爆炸之后的宇宙，宇宙怡然自得地，一边自我膨胀，一边给自己降温……
    \bitem
  \item[1]{{\blue 快乐大家庭:} \scriptsize 一开始，正负电子、光子、中微子等互相散射，处于热平衡（因此有共同的温度）。}
  \item[2]{{\blue 中微子不跟大家玩了:} \scriptsize 宇宙膨胀导致碰撞率快速下降，当温度降到大约 $\lesssim 1.5\mathrm{MeV}$ 时，正反中微子之间，以及中微子和正负电子之间基本不再有散射（中微子脱耦）。不过，中微子的温度仍然和光子-正负电子气体的温度保持一致（因两者都按照 $T\propto a^{-1}$ 演化）。}
  \item[3]{{\blue 正负电子挂了：} \scriptsize 不久之后，当温度降到电子的静止质量 ($\approx 0.5\mathrm{MeV}$)之下时，从光子产生正负电子也变得非常困难。绝大部分正电子和电子都湮没为光子，消失在历史的长河中……在这一步，正负电子把它们的\sout{功力}熵传递给了光子，导致光子温度骤然上升，超过了中微子的温度。}
  \item[4]{{\blue 一切又归于平静：} \scriptsize 在绝大部分正负电子湮灭之后，光子和中微子的温度比例就一直保持不变了，都是按 $T\propto a^{-1}$ 的规律变化。}    
    \eitem
  \end{frame}

  \begin{frame}
    \frametitle{这个故事的一些细节的说明}


    \bitem
    \item[1]{脱耦温度的精确计算有亿点点难（量子场论II，8课时走起），不过估算数量级并不难，我们之后会回顾这个问题。}
    \item[2]{脱耦的有三代中微子和反中微子；因为中微子不直接参与电磁相互作用，不存在“正反中微子湮灭为光子”的情况。(高能的正反中微子能湮灭为正负电子，但脱耦后后这个很难发生了）}
    \item[3]{脱耦的三代中微子和反中微子都只有一个内禀自由度。因为中微子都是左手征的，反中微子都是右手征的。}      
    \item[4]{因为正反物质不对称，正负电子湮灭之后还留下了（相对于湮灭前而言）极少量电子，后续上演了各种（可能再次导致你萌生退课念头的）好戏，但它们拥有的熵和光子相比有亿点低（字面意思），所以不会影响光子整体的平均温度。}
      \item[5]{中微子也可能存在$\sim \frac{\mu}{T}\lesssim 10^{-9}$ 量级的正反粒子不对称。我们在后续计算中将忽略这么微小的差异，丢掉化学势 $\mu$。}
        \eitem

  \end{frame}

  \begin{frame}

    为了计算正负电子湮灭后光子温度和中微子温度之比，我们来回顾一下玻色子（光子）气体和费米子（中微子）气体的熵——

    \addfig{1}{HEwarn.jpg}
  \end{frame}
  

  \begin{frame}
    \frametitle{玻色分布和费米分布：分布函数}
    你显然记得(\lfig{0.2}{tinyblackq.jpg})热统课里的（无势能，非极端简并的情况下）玻色分布和费米分布
    $$\frac{dn}{d p} = \frac{g}{2\pi^2}\frac{p^2}{e^{\frac{\sqrt{p^2+m^2}- \mu}{T}}\mp 1}$$
    这里的 $m$ 是静止质量，$\mu$ 是化学势, $p$ 是动量大小，$dn$ 是分布在 $p$ 和 $p+dp$ 间的粒子数密度, $g$ 是粒子内禀自由度个数（例如，光子的 $g=2$，因为它有两种螺旋度）。$-$对应玻色子，$+$对应费米子。


    也可以把它们写成单位能量区间的分布：
    $$\frac{dn}{d\epsilon} = \frac{g}{2\pi^2}\frac{\epsilon  \,\sqrt{\epsilon^2-m^2}}{e^{\frac{\epsilon- \mu}{T}}\mp 1} ,\ \ (\epsilon\ge m)$$


  \end{frame}


  \begin{frame}
    \frametitle{玻色分布和费米分布：粒子数密度}
    对 $d\epsilon$ 积分可以得到粒子数密度：
    $$ n = \frac{g}{2\pi^2}\int_m^\infty \frac{\epsilon\sqrt{\epsilon^2-m^2}d\epsilon}{e^{\frac{\epsilon- \mu}{T}}\mp 1}=   gK_{\mp}(1,\frac{m}{T},\frac{\mu}{T}) T^3.$$
    这里的函数
     $$K_{\mp}(n, \alpha, \beta) := \frac{1}{2\pi^2}\int_\alpha^\infty \frac{x^n\sqrt{x^2-\alpha^2}}{e^{x-\beta}\mp 1} dx $$
    (\question emm, 这个积分看起来有亿点点难）
  \end{frame}

  \begin{frame}
    \frametitle{玻色分布和费米分布：能量密度和压强}
    以 $\epsilon$ 为权重积分可以得到能量密度：
    $$ \rho = \frac{g}{2\pi^2}\int_m^\infty \frac{\epsilon^2\sqrt{\epsilon^2-m^2}d\epsilon}{e^{\frac{\epsilon- \mu}{T}}\mp 1}=  gK_{\mp}(2,\frac{m}{T},\frac{\mu}{T}) T^4.$$
    以 $\frac{p^2}{3\epsilon}$ 为权重积分可以得到压强（思考下为什么）：
    $$P = \frac{g}{3}\left[ K_{\mp}(2,\frac{m}{T},\frac{\mu}{T}) -  \left(\frac{m}{T}\right)^2 K_{\mp}(0,\frac{m}{T},\frac{\mu}{T})\right] T^4.$$

    \notes{对于零质量粒子，可以立即在上面的结果中验证 $P=\frac{1}{3}\rho$ }
  \end{frame}

  \begin{frame}
    \frametitle{玻色分布和费米分布：熵密度}
    根据自由焓（也叫Gibbs自由能）的定义：
    $$ G = U - TS + PV$$
    这里 $U$ 是内能，$V$ 是体积；以及化学势$\mu$是粒子平均自由焓，就有
    $$ n\mu = \rho - Ts + P$$
    这里 $\rho, s$ 分别是单位体积的能量密度和熵密度。这样可以解出：
    $$ s = g\left[\frac{4}{3}K_{\mp}(2,\frac{m}{T},\frac{\mu}{T}) - \frac{1}{3} \left(\frac{m}{T}\right)^2 K_{\mp}(0,\frac{m}{T},\frac{\mu}{T})-\frac{\mu}{T}K_{\mp}(1,\frac{m}{T},\frac{\mu}{T}) \right]T^3$$
  \end{frame}



  \begin{frame}
    \frametitle{
    平衡态玻色气体和费米气体的粒子数密度，能量密度，压强和熵密度的公式总结如下：}
    {\blue \small
      \bea
      n &=&  gK_{\mp}(1,\frac{m}{T},\frac{\mu}{T}) T^3; \newl
      \rho &=&  gK_{\mp}(2,\frac{m}{T},\frac{\mu}{T}) T^4; \newl     
      P &=& \frac{g}{3}\left[ K_{\mp}(2,\frac{m}{T},\frac{\mu}{T}) -  \left(\frac{m}{T}\right)^2 K_{\mp}(0,\frac{m}{T},\frac{\mu}{T})\right] T^4; \newl
      s &=& g\left[\frac{4}{3}K_{\mp}(2,\frac{m}{T},\frac{\mu}{T}) - \frac{1}{3} \left(\frac{m}{T}\right)^2 K_{\mp}(0,\frac{m}{T},\frac{\mu}{T})-\frac{\mu}{T}K_{\mp}(1,\frac{m}{T},\frac{\mu}{T}) \right]T^3.
      \eea
    }
    这里
    {\blue $$K_{\mp}(n, \alpha, \beta) = \frac{1}{2\pi^2}\int_\alpha^\infty \frac{x^n\sqrt{x^2-\alpha^2}}{e^{x-\beta}\mp 1} dx $$}
    约定玻色子取 $K_{-}$，费米子取 $K_+$。
    
  \end{frame}


  \begin{frame}
    \frametitle{别忙着点退课键}
    好消息是：就本课程的目的而言，我们只要会算 $m=\mu=0$ 的最简单情况就可以了。

    \addfig{1}{guakexian.jpg}

    (之所以完整地讲这些，是考虑到万一你要研究原初重子数生成等秃头课题呢……)

    
  \end{frame}
  

  \begin{frame}
    \frametitle{$K_{\mp}$的一些特殊值}
    用围道积分或者任何 \sout{令你头秃的}你喜欢的 技巧，可以很轻松地推出

    {\blue $$K_{-}(n,0,0) = \frac{(n+1)!}{2\pi^2}\zeta(n+2);\ \ K_{+}(n,0,0) = \left(1-\frac{1}{2^{n+1}}\right)K_{-}(n,0,0) $$ }   
    这里的
    $$\zeta(s) = \sum_{k=1}^\infty \frac{1}{k^s}$$
    是著名的 Riemann-zeta 函数(你在试图证明黎曼猜想时一定见过吧)。我们知道 $\zeta(2) = \frac{\pi^2}{6}$,  $\zeta(4) =\frac{\pi^4}{90}$，以及 $\zeta(3) = 1.202056903\ldots$，$\zeta(5) = 1.036927755\ldots$
  \end{frame}


  \begin{frame}
    \frametitle{极端相对论情形}
    在极端相对论并且无明显正反物质不对称的情形（$T\gg m,\mu$），可以近似取 $\frac{m}{T}=0, \frac{\mu}{T} = 0$，并利用刚才给出的 $K_{\mp}(n, 0, 0)$ 的值，得到:

    \skipline

        \bmini{0.48}
        {\scriptsize 极端相对论零化学势玻色子，}
    {\blue
      \bea
      n_{\rm boson} &=&  \frac{g\zeta(3)}{\pi^2} T^3; \newl
      \rho_{\rm boson} &=&  \frac{g\pi^2}{30} T^4; \newl     
      P_{\rm boson} &=& \frac{g\pi^2}{90}T^4; \newl
      s_{\rm boson} &=& \frac{2g\pi^2}{45}T^3.
      \eea
    }
      \emini
    \bmini{0.48}
     {\scriptsize 极端相对论零化学势费米子，}
    {\blue
      \bea
      n_{\rm fermion} &=& \frac{3}{4}n_{\rm boson}; \newl
      \rho_{\rm fermion} &=& \frac{7}{8}\rho_{\rm boson}; \newl
      P_{\rm fermion} &=& \frac{7}{8}P_{\rm boson}; \newl
      s_{\rm fermion} &=& \frac{7}{8} s_{\rm boson}.
      \eea
    }
      \emini
      

  \end{frame}
  


  \begin{frame}
    \frametitle{正负电子湮灭导致的升温}
    在 $T\gg \SIMeV$ 时，单位共动体积内的正负电子($g=2+2=4$)、光子($g=2$)的总熵是：
    $$ sa^3 = (4\times \frac{7}{8} + 2) \frac{2\pi^2}{45}(aT)^3 = \frac{11}{2}\frac{2\pi^2}{45}(aT)^3 $$
    根据宇宙学原理，共动体积之间互间等温且没有热量的净流动，因此共动体积内的熵守恒，所以这个阶段 $aT$ 保持不变。

    当温度降到 $0.5\SIMeV$ 附近，正负电子（湮灭）自由度消失：
    $$ sa^3 = 2 \frac{2\pi^2}{45}(aT)^3 $$    
    也就是说，这时的光子气体的 $aT$ 必须涨到之前的光子-正负电子气体的 $aT$ (这也是中微子的 $aT$) 的 $\left(\frac{11}{4}\right)^{1/3} \approx 1.4$ 倍，才能保持 $sa^3$ 不变。
  \end{frame}
  

  \begin{frame}
    \frametitle{宇宙中微子背景(CNB)的温度}

    All in all, 结论就是：在正负电子湮灭之后（$T\ll 0.5\SIMeV$） ，宇宙的光子背景温度是中微子背景温度的 $1.4$ 倍。
    
    我们今天看到的宇宙光子背景（宇宙微波背景辐射，CMB）的温度是 $T_{\rm CMB} = 2.726 \SIK$，由此可以推算出CNB的等效温度为：

    $$ T_{\rm CNB} = \left(\frac{4}{11}\right)^{1/3}T_{\rm CMB} \approx 1.95\SIK$$

    可惜的是：观测中微子背景及其困难，以目前的实验技术手段毫无希望检验这个预言。
  \end{frame}

  \begin{frame}
    \frametitle{$N_{\rm eff}=3.046$ 是什么鬼？}
    我们说脱耦留下的中微子有三代（$e$中微子，$\mu$子中微子，$\tau$子中微子），但在文献中经常看到 $N_{\rm eff} = 3.046$. 这是因为中微子在脱耦时并不完全处于热平衡态，有个小的修正。

    \skiplines
    
    所以，中微子的等效 $g = 2N_{\rm eff} = 6.092$。
  \end{frame}

  
  \begin{frame}
    \frametitle{$\Omega_r$的计算}
    如果中微子质量可以忽略，那么可以把它当成辐射形式能量。那么光子($g=2$)能量密度：
    $$\rho_{\gamma 0} = \frac{2\pi^2}{30}T_{\rm CMB}^4 $$
    和中微子($g=2N_{\rm eff}$)能量密度
    $$\rho_{\nu 0} = \frac{7}{8}\frac{2N_{\rm eff}\pi^2}{30}T_{\rm CNB}^4 = \frac{7}{8}\frac{2N_{\rm eff}\pi^2}{30}\left(\frac{4}{11}\right)^{4/3}T_{\rm CMB}^4  $$
    一共贡献了(Planck单位制下）：

      $$\Omega_r = \frac{8\pi}{3H_0^2} \frac{\pi^2\left(1+\frac{7}{8}\left(\frac{4}{11}\right)^{4/3}N_{\rm eff}\right)T_{\rm CMB}^4}{15}$$
  \end{frame}

  \begin{frame}
    \frametitle{好用的形式}
    把 $T_{\rm CMB}=2.726\SIK$, $N_{\rm eff}=3.046$ 代入并使用 units.py 进行普通单位制到Planck单位制的转化，得到： 
    $$\Omega_r  = \frac{2.4748\times(1+0.2271N_{\rm eff})}{h^2}\times 10^{-5}=\frac{4.187}{h^2}\times 10^{-5}$$
    这里的 $h$ 由 $H_0 = 100h\,\mathrm{km/s/Mpc}$ 定义。

    \skipline
    
    代入 $h\approx 0.7$ 可以得到 $\Omega_r$ 大约为 $8.5\times 10^{-5}$，其中光子贡献了约 $60\%$，中微子贡献了约 $40\%$。
  \end{frame}

  
  \section{Neutrinos in Late Universe}
  \secpage{晚期宇宙的中微子}{你以为高能预警已经结束了？不，并没有。}

  \begin{frame}
    \frametitle{中微子质量非零}
    前面我们假设了中微子质量可以忽略。实际上，中微子振荡实验已经确定三代中微子质量之和有个大约为 $0.06\SIeV$ 的下限（在质量反常排序模型中这个下限是 $0.1\SIeV$ 左右）。也就是说，在 $z\lesssim 100$ 的晚期宇宙，把中微子当成辐射形式能量不太靠谱。

    \skipline

    好在这对宇宙膨胀的计算只是个小的修正：因为在红移 $z\lesssim 300$，中微子的能量占比很小（除非中微子质量远超过 $0.1\SIeV$ 的量级，但这会导致很多和观测的不一致性）。
  \end{frame}
  


  \begin{frame}
    \frametitle{中微子的动量变化规律}
    为了修正以前的粗糙计算，我们希望在红移 $z$ 更加精确地计算中微子的能量密度和压强。

    \skipline
    
    不妨设某个中微子的四维动量是 $(p^t, p^r, 0, 0)$，测地线方程为：
    $$ \frac{dp^t}{dt} + \frac{a\dot a}{1-kr^2}\frac{(p^r)^2}{p^t} =0 $$
    再利用 $(p^t)^2 - \frac{a^2}{1-kr^2}(p^r)^2 = m^2$，上式等价于
    $ a\sqrt{(p^t)^2-m^2}$ 守恒。也就是说，{\blue FRW坐标系的静止观测者看到的自由粒子的三维动量大小都按 $\frac{1}{a}$ 衰减}。
  \end{frame}


  \begin{frame}
    \frametitle{中微子分布函数}
    由于 $q\equiv ap$ 是守恒的（这里 $p$ 是FRW坐标系的静止观测者看到的中微子三维动量大小），设脱耦时的尺度因子为 $a^*$，温度为 $T^*$，则每种中微子的分布函数都是：

    $$\left(\frac{a}{a^*}\right)^3dn = \frac{1}{2\pi^2}\frac{\left(\frac{q}{a^*}\right)^2d\frac{q}{a^*}}{e^{\frac{q}{a^*T^*}}+1}$$
    因为随着宇宙的膨胀，$a^3dn$ 是守恒量，所以我们把 $dn$ 乘以了 $\left(\frac{a}{a^*}\right)^3$ 以换算到脱耦时刻。

    注意到 $a^*T^* = T_{\rm CNB}=1.95\SIK$ (对质量不可忽略的中微子而言，$T_{\rm CNB}$ 只具有分布函数中的参数的意义)， 就有

    $$ a^3dn = \frac{1}{2\pi^2}\frac{q^2dq}{e^{\frac{q}{T_{\rm CNB}}}+1}$$

  \end{frame}
  

  \begin{frame}
    \frametitle{中微子能量密度}
    每种中微子的能量密度 $\rho$ 就满足
    $$\rho a^3 = \frac{1}{2\pi^2}\int_0^\infty \frac{q^2dq\sqrt{\frac{q^2}{a^2}+m^2}}{e^{\frac{q}{T_{\rm CNB}}}+1}$$    
    做变量替换 $x = q/T_{\rm CNB}$ 可以得到：
    $$\rho a^4 = \mathrm{I}_\rho\left(\frac{ma}{T_{\rm CNB}}\right) T^4_{\rm CNB} $$
    这里的
    $$ \mathrm{I}_\rho(\lambda) \equiv  \frac{1}{2\pi^2}\int_0^\infty \frac{x^2\sqrt{x^2+\lambda^2}}{e^x+1}dx$$    
  \end{frame}

  \begin{frame}
    \frametitle{$\mathrm{I}_\rho(\lambda)$ 的计算}
    当 $\lambda \ll 1$，有
    $$ \mathrm{I}_\rho(\lambda) \approx \frac{7\pi^2}{240} + \frac{1}{48}\lambda^2 $$
    当 $\lambda \gg 1$，有
    $$ \mathrm{I}_\rho(\lambda) \approx \frac{3\zeta(3)}{4\pi^2}\lambda + \frac{45\zeta(5)}{8\pi^2}\frac{1}{\lambda} $$ %$$\sum_{n=0}^\infty  \left(1-\frac{1}{2^{2n+2}}\right)\begin{pmatrix}\frac{1}{2} \\ n \end{pmatrix}\frac{(2n+2)!\zeta(2n+3)}{2\pi^2}\lambda^{1-2n}$$
  \end{frame}


  \begin{frame}
    \frametitle{中微子压强}
    每种中微子的压强 $p$ 就满足
    $$pa^3  = \frac{1}{2\pi^2}\int_0^\infty \frac{\frac{q^2}{a^2}}{3\sqrt{\frac{q^2}{a^2}+m^2}}\frac{q^2dq}{e^{\frac{q}{T_{\rm CNB}}}+1}$$    
    做变量替换 $x = q/T_{\rm CNB}$ 可以得到：
    $$ pa^4 = \mathrm{I}_p\left(\frac{ma}{T_{\rm CNB}}\right) T^4_{\rm CNB} $$
    这里的
    $$ \mathrm{I}_p(\lambda) \equiv  \frac{1}{6\pi^2}\int_0^\infty \frac{x^4}{\sqrt{x^2+\lambda^2}\left(e^x+1\right)}dx$$    
  \end{frame}

  \begin{frame}
    \frametitle{$\mathrm{I}_p(\lambda)$ 的计算}
    当 $\lambda \ll 1$，有
    $$ \mathrm{I}_p(\lambda) \approx \frac{7\pi^2}{720} - \frac{1}{144}\lambda^2 $$
    当 $\lambda \gg 1$，有
    $$ \mathrm{I}_p(\lambda) \approx \frac{15\zeta(5)}{4\pi^2}\frac{1}{\lambda} $$ %$$\sum_{n=0}^\infty  \left(1-\frac{1}{2^{2n+4}}\right)\begin{pmatrix}-\frac{1}{2} \\ n \end{pmatrix}\frac{(2n+4)!\zeta(2n+5)}{6\pi^2}\lambda^{-1-2n}$$

  \end{frame}


  \begin{frame}
    \frametitle{代码}
    网站上分享了简单的计算 $I_\rho$ 和 $I_p$，以及考虑到中微子质量非零后的宇宙年龄、各种距离、共动体积等的计算。

    \addfig{2.}{dudong.jpg}

    为了让你读得懂代码，我牺牲了一点点精度，中微子的质量修正的计算误差只控制在千分之一左右。不过，考虑到中微子质量修正只发生在晚期宇宙，这时中微子的能量占比不超过百分之几，所以总体计算（如$H(z)$)误差大致仍控制在 $10^{-5}$ 量级。
  \end{frame}
  
  
    \ech
\end{document}




  
